\documentclass[11pt]{article}

\usepackage[T1]{fontenc}
\usepackage[utf8]{inputenc}
\usepackage{lmodern}
\usepackage[margin=1in]{geometry}
\usepackage{microtype}
\usepackage{amsmath,amssymb,amsthm,mathtools}
\usepackage{booktabs,tabularx,longtable,array}
\usepackage{enumitem}
\usepackage{xcolor}
\usepackage{fancyhdr}
\usepackage{hyperref}
\usepackage[nameinlink,capitalise,noabbrev]{cleveref}

\definecolor{linkblue}{RGB}{24,72,124}
\definecolor{citegreen}{RGB}{23,102,69}
\definecolor{shadegray}{RGB}{246,247,249}
\hypersetup{
  colorlinks=true,
  linkcolor=linkblue,
  citecolor=citegreen,
  urlcolor=linkblue,
  pdftitle={Exact q-Binomial Extraction of Rvachev Up-Function Values from Finite Sinc-Product Splines},
  pdfauthor={Research article prepared for the ProveIt FabiusFunction project},
  pdfsubject={Dyadic spline values, finite local deconvolution, geometric progressions, q-binomial extrapolation},
  pdfkeywords={Rvachev up-function, Fabius function, dyadic rational, sinc product, spline, q-binomial, q-Pochhammer, Richardson extrapolation, Thue-Morse}
}

\pagestyle{fancy}
\fancyhf{}
\fancyhead[L]{\small Exact dyadic extraction from finite sinc products}
\fancyhead[R]{\small Rvachev up-function}
\fancyfoot[C]{\thepage}
\setlength{\headheight}{14pt}

\newtheorem{theorem}{Theorem}[section]
\newtheorem{lemma}[theorem]{Lemma}
\newtheorem{proposition}[theorem]{Proposition}
\newtheorem{corollary}[theorem]{Corollary}
\newtheorem{definition}[theorem]{Definition}
\newtheorem{remark}[theorem]{Remark}
\newtheorem{example}[theorem]{Example}

\newcommand{\R}{\mathbb{R}}
\newcommand{\Z}{\mathbb{Z}}
\newcommand{\N}{\mathbb{N}}
\newcommand{\dd}{\,\mathrm{d}}
\newcommand{\upf}{\operatorname{up}}
\newcommand{\sinc}{\operatorname{sinc}}
\newcommand{\e}{\mathrm{e}}
\newcommand{\ii}{\mathrm{i}}
\newcommand{\ind}{\mathbf{1}}
\newcommand{\pos}[1]{\left(#1\right)_{+}}
\newcommand{\qbinom}[3]{\genfrac{[}{]}{0pt}{}{#1}{#2}_{#3}}
\newcommand{\Qpoch}[2]{\left(#1;#1\right)_{#2}}
\newcommand{\Shift}{\mathsf{E}}

\title{\bfseries Exact \(q\)-Binomial Extraction of Rvachev Up-Function Values\\
from Finite Sinc-Product Splines}
\author{Research article prepared for the \texttt{ProveIt/FabiusFunction} project}
\date{2 September 2026}

\begin{document}
\maketitle

\begin{abstract}
Let
\[
 \widehat{\upf}(t)=\prod_{k=0}^{\infty}\sinc\!\left(\frac{\pi t}{2^k}\right),
 \qquad \sinc z=\frac{\sin z}{z},
\]
and let \(\upf_n\) be the inverse Fourier transform of the prefix product with
\(0\le k\le n\).  Fix a dyadic rational \(x\in[-1,1]\), and write
\(d=d(x)=\min\{r\ge0:2^r x\in\Z\}\) and \(M=\lfloor d/2\rfloor\).
The principal result is the exact eventual identity
\[
 \upf_n(x)=\sum_{m=0}^{M}(-1)^m a_m\,
 \upf^{(2m)}(x)\,4^{-m(n+1)},\qquad n\ge d,
\]
where the rational coefficients \(a_m\) are defined by
\[
 \frac{1}{\widehat{\upf}(z)}
 =\sum_{m\ge0}a_m(2\pi z)^{2m}.
\]
Thus the sequence \(q_n=\upf_n(x)\) is, after at most \(d\) initial levels,
exactly a constant plus finitely many signed geometric progressions with ratios
\(4^{-1},4^{-2},\ldots,4^{-M}\).  The proof is not an appeal to a merely
asymptotic Fourier expansion.  It uses two finite mechanisms: the omitted tail
is a scaled copy of \(\upf\), and for \(n\ge d\) its support lies exactly in one
polynomial spline cell centered at \(x\); moreover, every derivative
\(\upf^{(r)}(x)\) with \(r>d\) vanishes.

Putting \(Q=1/4\), geometric Lagrange interpolation at the nodes
\(1,Q,\ldots,Q^M\) gives the requested one-line extraction formula
\[
 \boxed{\displaystyle
 \upf(x)=\frac{1}{(Q;Q)_M}
 \sum_{j=0}^{M}(-1)^{M-j}Q^{\binom{M-j+1}{2}}
 \qbinom{M}{j}{Q}\,\upf_{n+j}(x),\qquad n\ge d.}
\]
All quantities on the right are rational.  Hence this formula simultaneously
extracts the limit exactly and gives a direct finite proof of the rationality of
\(\upf(x)\) at dyadic points.  We also give the complete \(q\)-Vandermonde
system for all geometric coefficients, an explicit inverse formula, a finite
\(q\)-binomial annihilating recurrence, stability bounds, worked rational
examples, and exact-arithmetic Python verification.
\end{abstract}

\tableofcontents
\clearpage

\section{Introduction}

The Rvachev up-function is a compactly supported, even, infinitely
differentiable function whose Fourier transform is an infinite dyadic sinc
product.  It may equally be viewed as the probability density of an infinite
sum of independent centered uniform random variables.  Truncating the Fourier
product produces a compactly supported box spline.  The resulting finite
approximants are especially rigid: their knots form a uniform dyadic grid and
the jumps of the top derivative carry the Thue--Morse signs.

The documentation in the \texttt{ProveIt} repository develops these finite
sinc products, their exact truncated-power representation, the self-similar
head--tail factorization, a uniform all-orders differential expansion, and
geometric Richardson filters with closed Gaussian \(q\)-binomial weights
\cite{ReshetnikovRepo}.  Those ingredients suggest the following sharper
pointwise question.

\begin{quote}
For a fixed dyadic rational \(x\), does the asymptotic series in powers of
\(4^{-n}\) become an \emph{exact finite formula}, and can the true value
\(\upf(x)\) be recovered from finitely many rational spline values by one
closed expression?
\end{quote}

The answer is yes.  The decisive observation is geometric.  Once the spline
level reaches the denominator depth of \(x\), the point \(x\) is the midpoint
of a spline cell, and the entire support of the omitted scaled tail is exactly
that cell.  Convolution by the tail therefore acts on one finite-dimensional
polynomial space.  Its inverse is not an infinite pseudodifferential series on
that space: it is a finite nilpotent differential operator.  The coefficients
of this inverse are precisely the reciprocal-product coefficients already
visible in Fourier analysis.  A second finite phenomenon then occurs: the
Taylor jet of \(\upf\) at a dyadic point terminates.  Combining the two gives
an exact sum of geometric modes.

The limit is the constant coefficient of a polynomial sampled at the geometric
nodes \(1,Q,\ldots,Q^M\), where \(Q=1/4\).  Evaluating that polynomial at zero
is elementary Lagrange interpolation, and simplification of the Lagrange
weights produces the advertised \(q\)-Pochhammer/Gaussian-binomial formula.
This is also a particularly well-conditioned extrapolation: the total
variation of the weights stays below \(1.97\), uniformly in the dyadic depth.

\subsection*{Main conclusions}

For a dyadic point of depth \(d\), set \(M=\lfloor d/2\rfloor\).  Then:
\begin{enumerate}[leftmargin=2.2em]
\item the guaranteed transient ends at level \(n=d\);
\item \(q_n=\upf_n(x)\) is exactly a polynomial of degree at most \(M\) in
      \(Q^{n+1}\), and the degree is \(M\) for every nonintegral reduced
      dyadic point;
\item \(M+1\) consecutive rational values determine every mode coefficient;
\item the constant term \(\upf(x)\) is obtained by a single universal
      \(q\)-binomial filter depending only on \(d\), not on the numerator of
      \(x\);
\item the same sequence satisfies an exact order-\(M+1\) recurrence whose
      coefficients are Gaussian \(q\)-binomials.
\end{enumerate}

All numerical calculations in the accompanying script use
\texttt{fractions.Fraction}.  They are verification and illustration only;
every mathematical assertion used by the extraction method is proved below.

\section{Fourier normalization and the finite box splines}

\subsection{The infinite product}

We use the Fourier convention
\begin{equation}\label{eq:fourier}
 \widehat f(t)=\int_{\R}f(x)\e^{-2\pi\ii tx}\dd x,
 \qquad
 f(x)=\int_{\R}\widehat f(t)\e^{2\pi\ii tx}\dd t.
\end{equation}
For \(a>0\), let
\[
 b_a(x)=\frac{1}{2a}\ind_{[-a,a]}(x).
\]
Then
\[
 \widehat b_a(t)=\sinc(2\pi at).
\]
Consequently the normalized up-function is the probability density with
Fourier transform
\begin{equation}\label{eq:Phi}
 \Phi(t):=\widehat{\upf}(t)
 =\prod_{k=0}^{\infty}\sinc\!\left(\frac{\pi t}{2^k}\right),
\end{equation}
and support \([-1,1]\).  This is the normalization used in the classical work
of Rvachev and Arias de Reyna \cite{Rvachev1990,Arias1982}.

For \(n\ge0\), define the user-indexed prefix
\begin{equation}\label{eq:prefix-user}
 \Phi_n(t)=\prod_{k=0}^{n}\sinc\!\left(\frac{\pi t}{2^k}\right),
 \qquad \widehat{\upf_n}=\Phi_n.
\end{equation}
There are \(n+1\) factors, so
\begin{equation}\label{eq:prefix-convolution}
 \upf_n=b_{2^{-1}}*b_{2^{-2}}*\cdots*b_{2^{-(n+1)}}.
\end{equation}
In particular, \(\upf_n\) is a probability density and is a polynomial spline
of degree \(n\).

\subsection{Random-sum and exact tail factorization}

Let \(U_0,U_1,\ldots\) be independent random variables, each uniform on
\([-1,1]\), and put
\begin{equation}\label{eq:random-sums}
 S_n=\sum_{k=0}^{n}2^{-k-1}U_k,
 \qquad
 X=\sum_{k=0}^{\infty}2^{-k-1}U_k.
\end{equation}
Then \(\upf_n\) and \(\upf\) are the densities of \(S_n\) and \(X\),
respectively.  The omitted tail is
\[
 \sum_{k=n+1}^{\infty}2^{-k-1}U_k
 =2^{-(n+1)}\sum_{j=0}^{\infty}2^{-j-1}U_{n+1+j}.
\]
Thus, with
\begin{equation}\label{eq:delta}
 \delta_n=2^{-(n+1)},
\end{equation}
we have the distributional identity
\begin{equation}\label{eq:head-tail-random}
 X\overset{d}=S_n+\delta_n X',
\end{equation}
where \(X'\) is an independent copy of \(X\).  At the density and Fourier
levels this becomes
\begin{equation}\label{eq:head-tail}
 \Phi(t)=\Phi_n(t)\Phi(\delta_n t),
 \qquad
 \upf=\upf_n*\rho_{\delta_n},
 \qquad
 \rho_\delta(y)=\delta^{-1}\upf(y/\delta).
\end{equation}
The support of \(\rho_\delta\) is \([-\delta,\delta]\).

\section{Exact spline arithmetic and dyadic cell geometry}

\subsection{Thue--Morse truncated-power formula}

Let
\[
 \tau_k=(-1)^{s_2(k)},
\]
where \(s_2(k)\) is the sum of the binary digits of \(k\).  Repeatedly
differentiating the boxes in \cref{eq:prefix-convolution}, or integrating the
resulting atomic top derivative, gives the following exact formula.

\begin{proposition}[Exact finite-spline formula]\label{prop:truncated-power}
For every \(n\ge0\), away from the knots when \(n=0\),
\begin{equation}\label{eq:truncated-power}
 \boxed{\displaystyle
 \upf_n(x)=\frac{2^{n(n+1)/2}}{n!}
 \sum_{k=0}^{2^{n+1}-1}\tau_k
 \pos{x+1-2^{-(n+1)}-\frac{k}{2^n}}^{n}.}
\end{equation}
Its knots are
\begin{equation}\label{eq:knots}
 t_{n,k}=-1+2^{-(n+1)}+\frac{k}{2^n}
        =-1+(2k+1)2^{-(n+1)},
 \qquad 0\le k<2^{n+1}.
\end{equation}
The support is \([-1+2^{-(n+1)},1-2^{-(n+1)}]\), and on each open knot cell
\(\upf_n\) is a polynomial of degree at most \(n\).
\end{proposition}

\begin{proof}
For a centered box,
\[
 Db_a=\frac{1}{2a}(\delta_{-a}-\delta_a).
\]
Taking one derivative from every factor in \cref{eq:prefix-convolution} shows
that \(D^{n+1}\upf_n\) is a signed sum of \(2^{n+1}\) point masses.  Choosing
right endpoints according to a binary vector gives the sign \(\tau_k\), the
uniform location in \cref{eq:knots}, and the normalization
\[
 \prod_{j=0}^{n}2^j=2^{n(n+1)/2}.
\]
Integrating \(n+1\) times from \(-\infty\), using compact support, yields
\cref{eq:truncated-power}.  The support and polynomial-cell statements are
immediate.
\end{proof}

\begin{corollary}[Rational prefix values]\label{cor:prefix-rational}
If \(x\in\mathbb{Q}\), then \(\upf_n(x)\in\mathbb{Q}\).  In particular every
term \(q_n=\upf_n(x)\) is rational when \(x\) is dyadic.
\end{corollary}

\subsection{Dyadic depth and the centered-cell lemma}

\begin{definition}[Dyadic depth]\label{def:depth}
For a dyadic rational \(x\), define
\[
 d(x)=\min\{d\ge0:2^d x\in\Z\}.
\]
If \(x\notin\Z\), this says that \(x=a/2^d\) in lowest terms with \(a\) odd.
\end{definition}

The next elementary observation is the geometric reason for the exact
termination theorem.

\begin{lemma}[A dyadic point becomes a cell midpoint]\label{lem:centered-cell}
Let \(x\in(-1,1)\) be dyadic of depth \(d\).  For every \(n\ge d\), the two
knots adjacent to \(x\) in the spline \(\upf_n\) are
\[
 x-\delta_n\quad\text{and}\quad x+\delta_n,
 \qquad \delta_n=2^{-(n+1)}.
\]
Hence \((x-\delta_n,x+\delta_n)\) is one polynomial cell of \(\upf_n\), centered
at \(x\), and it is exactly as wide as the support of the omitted tail kernel
in \cref{eq:head-tail}.
\end{lemma}

\begin{proof}
Since \(n\ge d\), the number
\[
 J=2^n(x+1)
\]
is an integer.  Therefore
\[
 x=-1+\frac{J}{2^n}=-1+2J\,2^{-(n+1)}.
\]
By \cref{eq:knots}, the adjacent odd-grid points are
\[
 -1+(2J-1)2^{-(n+1)}=x-\delta_n,
 \qquad
 -1+(2J+1)2^{-(n+1)}=x+\delta_n.
\]
There is no knot strictly between them.
\end{proof}

\begin{remark}
The guaranteed threshold in the user's indexing is therefore exactly
\(n=d\).  Some special points may satisfy the final formula at earlier levels,
but no earlier universal threshold follows from the cell geometry.
\end{remark}

\section{Reciprocal-product coefficients and finite local deconvolution}

\subsection{Moments and the reciprocal sinc product}

Let
\begin{equation}\label{eq:moments}
 \mu_r=\int_{-1}^{1}t^r\upf(t)\dd t=\mathbb{E}[X^r].
\end{equation}
The density is even, so \(\mu_{2r+1}=0\).  Its moment-generating series is
\begin{equation}\label{eq:moment-series}
 B(s)=\mathbb{E}[\e^{sX}]
     =\sum_{m=0}^{\infty}\frac{\mu_{2m}}{(2m)!}s^{2m}.
\end{equation}
Because \(\Phi(z)=B(-2\pi\ii z)\), define \(a_m\) near the origin by
\begin{equation}\label{eq:reciprocal-def}
 \frac{1}{\Phi(z)}
 =\sum_{m=0}^{\infty}a_m(2\pi z)^{2m}.
\end{equation}
Substituting \(z=\ii s/(2\pi)\) gives the inverse moment series
\begin{equation}\label{eq:inverse-moment-series}
 \frac{1}{B(s)}=\sum_{m=0}^{\infty}(-1)^m a_m s^{2m}.
\end{equation}
Equivalently, if \(b_m=\mu_{2m}/(2m)!\), then
\begin{equation}\label{eq:convolution-inverse-coeff}
 \sum_{j=0}^{m}b_j(-1)^{m-j}a_{m-j}=\delta_{m0}.
\end{equation}

The logarithm of the sinc product gives an explicit rational description.  Put
\(Q=1/4\).  Then
\begin{equation}\label{eq:alpha}
 \frac1{\Phi(z)}
 =\exp\!\left(\sum_{j=1}^{\infty}\alpha_j(2\pi z)^{2j}\right),
 \qquad
 \alpha_j=\frac{|B_{2j}|}{2j(2j)!(1-Q^j)},
\end{equation}
where \(B_{2j}\) are the Bernoulli numbers.  Hence
\begin{equation}\label{eq:a-recurrence}
 a_0=1,
 \qquad
 m a_m=\sum_{j=1}^{m}j\alpha_j a_{m-j},
\end{equation}
and, using the complete exponential Bell polynomial \(Y_m\),
\begin{equation}\label{eq:a-bell}
 a_m=\frac1{m!}Y_m(1!\alpha_1,2!\alpha_2,\ldots,m!\alpha_m).
\end{equation}
In particular, every \(a_m\) is a positive rational number.  The first values
are
\begin{equation}\label{eq:first-a}
 a_0=1,
 \quad a_1=\frac1{18},
 \quad a_2=\frac{31}{16200},
 \quad a_3=\frac{2347}{42865200},
 \quad a_4=\frac{1904369}{1311675120000}.
\end{equation}

\subsection{An exact local inverse, not an asymptotic one}

For an arbitrary point, formally dividing by \(\Phi(\delta_n t)\) leads to an
all-orders asymptotic differential expansion.  At the centered dyadic cell the
same calculation is exact because it takes place on a finite-dimensional
polynomial space.

\begin{theorem}[Finite local deconvolution]\label{thm:local-deconvolution}
Let \(x\in(-1,1)\) be dyadic of depth \(d\), let \(n\ge d\), and let
\(0\le r\le n\).  Then
\begin{equation}\label{eq:exact-local-inverse}
 \boxed{\displaystyle
 \upf_n^{(r)}(x)=
 \sum_{m=0}^{\lfloor(n-r)/2\rfloor}
 (-1)^m a_m\,\delta_n^{2m}\,\upf^{(r+2m)}(x),
 \qquad \delta_n=2^{-(n+1)}.}
\end{equation}
This identity is exact for each such \(n\); there is no remainder term.
\end{theorem}

\begin{proof}
By \cref{lem:centered-cell}, \(\upf_n\) agrees on
\((x-\delta_n,x+\delta_n)\) with one polynomial \(g\) of degree at most \(n\).
From the exact convolution identity \cref{eq:head-tail}, differentiation in the
sense of distributions followed by evaluation at \(x\) gives, for
\(0\le r\le n\),
\begin{align}
 \upf^{(r)}(x)
 &=\int_{-1}^{1}\upf_n^{(r)}(x-\delta_n t)\upf(t)\dd t \notag\\
 &=\int_{-1}^{1}g^{(r)}(x-\delta_n t)\upf(t)\dd t.\label{eq:local-convolution-jet}
\end{align}
The possible cell endpoints correspond to \(t=\pm1\) and have measure zero.
Since \(g^{(r)}\) is a polynomial, its Taylor expansion is finite.  Odd moments
vanish, so
\begin{equation}\label{eq:triangular-forward}
 \upf^{(r)}(x)=
 \sum_{j=0}^{\lfloor(n-r)/2\rfloor}
 b_j\delta_n^{2j}\upf_n^{(r+2j)}(x),
 \qquad b_j=\frac{\mu_{2j}}{(2j)!}.
\end{equation}
This is a unit upper-triangular relation between the even subjets.  Its inverse
is determined by the reciprocal power series of \(\sum b_j s^{2j}=B(s)\).
By \cref{eq:inverse-moment-series}, the inverse coefficients are
\((-1)^m a_m\).  Because the derivative shift is nilpotent on polynomials of
degree at most \(n\), multiplying the two finite triangular sums uses only the
finite coefficient identities \cref{eq:convolution-inverse-coeff}; no
convergence argument is involved.  Applying the inverse relation to
\cref{eq:triangular-forward} proves \cref{eq:exact-local-inverse}.
\end{proof}

\begin{remark}[Operator interpretation]
On the local polynomial \(g\), convolution by the tail is the finite operator
\[
 B(-\delta_n D)g=B(\delta_n D)g,
\]
because \(B\) is even.  Its inverse is
\[
 B(\delta_n D)^{-1}
 =\sum_{m\ge0}(-1)^m a_m\delta_n^{2m}D^{2m},
\]
and the sum terminates automatically since \(D^{n+1}g=0\).  This is the local,
finite-dimensional meaning of division by the omitted Fourier factor.
\end{remark}

\section{Termination of the dyadic Taylor jet}

We next prove the second finite mechanism: sufficiently high derivatives of
\(\upf\) vanish at every dyadic point.

\subsection{Functional-differential equation}

The product identity
\[
 \Phi(t)=\sinc(\pi t)\Phi(t/2)
\]
is equivalent, under \cref{eq:fourier}, to
\begin{equation}\label{eq:up-diffeq}
 \upf'(x)=2\upf(2x+1)-2\upf(2x-1).
\end{equation}
Iterating this relation gives the standard Thue--Morse derivative formula.

\begin{lemma}[Iterated derivative identity]\label{lem:derivative-formula}
For every integer \(r\ge0\),
\begin{equation}\label{eq:derivative-formula}
 \boxed{\displaystyle
 \upf^{(r)}(x)=2^{r(r+1)/2}
 \sum_{k=0}^{2^r-1}\tau_k\,
 \upf\!\left(2^r x+2^r-1-2k\right).}
\end{equation}
\end{lemma}

\begin{proof}
The case \(r=0\) is tautological, and the case \(r=1\) is
\cref{eq:up-diffeq}.  Suppose the formula holds for \(r\).  Differentiating
introduces a factor \(2^r\), and applying \cref{eq:up-diffeq} to each summand
introduces another factor \(2\) and two children whose signs are related by
\[
 \tau_{2k}=\tau_k,
 \qquad
 \tau_{2k+1}=-\tau_k.
\]
The total new factor is \(2^{r+1}\), changing
\(2^{r(r+1)/2}\) into \(2^{(r+1)(r+2)/2}\), and the children enumerate
\(0\le k<2^{r+1}\).  This proves the induction step.
\end{proof}

\subsection{Dyadic flatness}

\begin{theorem}[Dyadic jet termination]\label{thm:dyadic-flatness}
Let \(x\in[-1,1]\) be dyadic of depth \(d\).  Then
\begin{equation}\label{eq:dyadic-vanishing}
 \upf^{(r)}(x)=0\qquad\text{for every }r>d.
\end{equation}
For a nonintegral reduced dyadic point, the highest even derivative permitted by
this bound is nonzero: if \(M=\lfloor d/2\rfloor\), then
\begin{equation}\label{eq:top-even-nonzero}
 \upf^{(2M)}(x)\ne0.
\end{equation}
\end{theorem}

\begin{proof}
First take \(r>d\).  Then \(2^r x\) is an even integer.  Every shift
\(2^r-1-2k\) in \cref{eq:derivative-formula} is odd, so every argument of
\(\upf\) in that formula is an odd integer.  Since \(\upf\) is supported on
\([-1,1]\) and is smooth, it vanishes at \(\pm1\) and outside the support.
Thus every summand is zero, proving \cref{eq:dyadic-vanishing}.  The same
argument includes \(x=0,\pm1\).

It remains to show \cref{eq:top-even-nonzero} for a nonintegral point
\(x=a/2^d\) in lowest terms.  We first record the elementary values
\begin{equation}\label{eq:up-special-values}
 \upf(0)=1,
 \qquad
 \upf(1/2)=\upf(-1/2)=1/2.
\end{equation}
Indeed, for \(0\le y\le1\), differentiating
\(h(y)=\upf(y)+\upf(y-1)\) and using \cref{eq:up-diffeq} shows that \(h\) is
constant; its integral over \([0,1]\) is the total mass \(1\), so \(h\equiv1\).
At \(y=0\), endpoint flatness gives \(\upf(0)=1\); at \(y=1/2\), evenness gives
\(2\upf(1/2)=1\).

If \(d=2M\) is even, put \(r=d\) in \cref{eq:derivative-formula}.  The arguments
are even integers spaced by two, and exactly one is zero.  Therefore
\[
 \upf^{(d)}(x)=\pm2^{d(d+1)/2}\upf(0)\ne0.
\]
If \(d=2M+1\) is odd, put \(r=d-1=2M\).  The arguments are half-integers
spaced by two, and exactly one of them is \(1/2\) or \(-1/2\); all others lie
outside \([-1,1]\).  Hence
\[
 \upf^{(d-1)}(x)=\pm2^{(d-1)d/2}\upf(1/2)\ne0.
\]
This is precisely \cref{eq:top-even-nonzero}.
\end{proof}

\begin{corollary}[Terminating Taylor polynomial]\label{cor:taylor-terminates}
At a dyadic point of depth \(d\), the formal Taylor series of \(\upf\) contains
no term of degree greater than \(d\).  This finite Taylor polynomial need not
represent \(\upf\) on any neighborhood; indeed the up-function is nowhere
analytic on its support \cite{Arias1982}.
\end{corollary}

\section{The exact eventual geometric formula}

We can now combine the centered-cell inverse with dyadic jet termination.

\begin{theorem}[Exact finite geometric expansion]\label{thm:main-geometric}
Fix a dyadic rational \(x\in[-1,1]\), let
\[
 d=d(x),\qquad M=\left\lfloor\frac d2\right\rfloor,
 \qquad Q=\frac14,
\]
and set \(q_n=\upf_n(x)\).  Then, for every \(n\ge d\),
\begin{equation}\label{eq:main-geometric}
 \boxed{\displaystyle
 q_n=\sum_{m=0}^{M}(-1)^m a_m\upf^{(2m)}(x)Q^{m(n+1)}.}
\end{equation}
Equivalently,
\begin{equation}\label{eq:constant-plus-modes}
 q_n=\upf(x)+\sum_{m=1}^{M}C_m(x)Q^{m(n+1)},
 \qquad
 C_m(x)=(-1)^m a_m\upf^{(2m)}(x).
\end{equation}
Thus the nonconstant terms are signed geometric progressions whose successive
ratios are
\[
 Q^m=4^{-m},\qquad 1\le m\le M.
\]
For a nonintegral reduced dyadic point, the coefficient of the highest mode
\(Q^{M(n+1)}\) is nonzero, so the geometric polynomial has degree exactly
\(M\).
\end{theorem}

\begin{proof}
For \(x\in(-1,1)\) and \(n\ge d\), apply
\cref{thm:local-deconvolution} with \(r=0\).  Its finite sum initially runs to
\(\lfloor n/2\rfloor\).  By \cref{thm:dyadic-flatness}, every derivative with
order greater than \(d\) vanishes, so only \(2m\le d\), that is
\(m\le M\), survive.  Since \(\delta_n^{2m}=4^{-m(n+1)}=Q^{m(n+1)}\), this is
\cref{eq:main-geometric}.  The endpoint sequences are identically zero, and
the central point sequence is identically one, so the same statement holds at
\(x=-1,0,1\).  Finally, \(a_M>0\) and
\(\upf^{(2M)}(x)\ne0\) by \cref{thm:dyadic-flatness}, proving exact degree for
nonintegral reduced points.
\end{proof}

\begin{remark}[Why the formula is stronger than the global expansion]
Fourier division gives, for fixed truncation order \(K\), an asymptotic formula
with a remainder \(O(4^{-K n})\) uniformly in \(x\).  At a dyadic point, the
present proof does not attempt to make that remainder small.  It removes it:
local convolution is inverted exactly on one spline polynomial, and all
remaining derivatives above the dyadic depth are exactly zero.
\end{remark}

\begin{corollary}[Exact rationality from finite splines]\label{cor:rationality}
For every dyadic \(x\in[-1,1]\), \(\upf(x)\in\mathbb{Q}\).  Moreover all mode
coefficients \(C_m(x)\) are rational.
\end{corollary}

\begin{proof}
The next section expresses \(\upf(x)\) and every \(C_m(x)\) as rational linear
combinations of finitely many rational prefix values from
\cref{cor:prefix-rational}.  This gives a finite proof independent of taking a
limit of rationals.  Rationality of all dyadic derivatives then also follows
from \cref{eq:derivative-formula}.
\end{proof}

\section{Geometric interpolation and the one-line extraction formula}

\subsection{The \(q\)-Vandermonde system}

Fix a starting level \(n_0\ge d\).  Rewrite \cref{eq:constant-plus-modes} as
\begin{equation}\label{eq:scaled-coefficients}
 q_{n_0+j}=\sum_{m=0}^{M}B_m Q^{jm},
 \qquad
 B_m=C_m(x)Q^{m(n_0+1)},
 \qquad 0\le j\le M.
\end{equation}
Thus \((B_0,\ldots,B_M)^T\) is the solution of the exact system
\begin{equation}\label{eq:q-vandermonde-system}
 \begin{pmatrix}
 1&1&1&\cdots&1\\
 1&Q&Q^2&\cdots&Q^M\\
 1&Q^2&Q^4&\cdots&Q^{2M}\\
 \vdots&\vdots&\vdots&&\vdots\\
 1&Q^M&Q^{2M}&\cdots&Q^{M^2}
 \end{pmatrix}
 \begin{pmatrix}B_0\\B_1\\\vdots\\B_M\end{pmatrix}
 =
 \begin{pmatrix}q_{n_0}\\q_{n_0+1}\\\vdots\\q_{n_0+M}\end{pmatrix}.
\end{equation}
Its determinant is
\begin{equation}\label{eq:q-vandermonde-det}
 \prod_{0\le r<s\le M}(Q^s-Q^r)
 =(-1)^{\binom{M+1}{2}}Q^{\binom{M+1}{3}}
  \prod_{s=1}^{M}(Q;Q)_s,
\end{equation}
which is nonzero because \(0<Q<1\).  Hence \(M+1\) consecutive values uniquely
determine every geometric coefficient.  The desired limit is simply
\[
 B_0=C_0(x)=\upf(x).
\]

\subsection{Lagrange evaluation at zero}

Let
\[
 P(z)=\sum_{m=0}^{M}B_mz^m.
\]
Then \(P(Q^j)=q_{n_0+j}\), and \(\upf(x)=P(0)\).  Lagrange interpolation gives
\begin{equation}\label{eq:lagrange-zero}
 P(0)=\sum_{j=0}^{M}W_{M,j}q_{n_0+j},
 \qquad
 W_{M,j}=\prod_{\substack{0\le\ell\le M\\\ell\ne j}}
 \frac{-Q^\ell}{Q^j-Q^\ell}.
\end{equation}
We now simplify the product.

For \(\ell<j\),
\[
 \prod_{\ell=0}^{j-1}(Q^j-Q^\ell)
 =(-1)^jQ^{j(j-1)/2}(Q;Q)_j,
\]
while for \(\ell>j\),
\[
 \prod_{\ell=j+1}^{M}(Q^j-Q^\ell)
 =Q^{j(M-j)}(Q;Q)_{M-j}.
\]
The numerator in \cref{eq:lagrange-zero} is
\((-1)^M Q^{M(M+1)/2-j}\).  Cancelling powers and signs yields
\begin{equation}\label{eq:weight-product}
 W_{M,j}=(-1)^{M-j}
 \frac{Q^{\binom{M-j+1}{2}}}{(Q;Q)_j(Q;Q)_{M-j}}.
\end{equation}
Using
\begin{equation}\label{eq:qbinom-def}
 \qbinom{M}{j}{Q}
 =\frac{(Q;Q)_M}{(Q;Q)_j(Q;Q)_{M-j}},
\end{equation}
we obtain the central extraction theorem.

\begin{theorem}[Exact \(q\)-binomial extraction]\label{thm:q-extraction}
Let \(x\in[-1,1]\) be dyadic of depth \(d\), put
\(M=\lfloor d/2\rfloor\) and \(Q=1/4\), and choose any \(n\ge d\).  Then
\begin{equation}\label{eq:q-extraction}
 \boxed{\displaystyle
 \upf(x)=
 \sum_{j=0}^{M}(-1)^{M-j}
 \frac{Q^{\binom{M-j+1}{2}}}{(Q;Q)_j(Q;Q)_{M-j}}\,q_{n+j}.}
\end{equation}
Equivalently, in the requested single Gaussian-binomial/\(q\)-Pochhammer form,
\begin{equation}\label{eq:q-extraction-gaussian}
 \boxed{\displaystyle
 \upf(x)=\frac{1}{(Q;Q)_M}
 \sum_{j=0}^{M}(-1)^{M-j}Q^{\binom{M-j+1}{2}}
 \qbinom{M}{j}{Q}\,\upf_{n+j}(x),
 \qquad Q=\frac14,\quad n\ge d.}
\end{equation}
\end{theorem}

\begin{proof}
The exact geometric formula makes \(q_{n+j}\) the values of a polynomial of
degree at most \(M\) at \(Q^j\).  Lagrange interpolation at zero proves
\cref{eq:lagrange-zero}, and the preceding product calculation proves
\cref{eq:q-extraction,eq:q-extraction-gaussian}.
\end{proof}

\subsection{Shift-operator form}

Let \(\Shift q_n=q_{n+1}\).  Every nonconstant mode \(Q^{m(n+1)}\) is an
eigenvector of \(\Shift\) with eigenvalue \(Q^m\), while the constant mode has
eigenvalue \(1\).  Therefore
\begin{equation}\label{eq:shift-filter}
 \boxed{\displaystyle
 \upf(x)=
 \frac{(\Shift-Q)(\Shift-Q^2)\cdots(\Shift-Q^M)}
      {(1-Q)(1-Q^2)\cdots(1-Q^M)}\,q_n,
 \qquad n\ge d.}
\end{equation}
The denominator is \((Q;Q)_M\).  Expanding the numerator by the finite
\(q\)-binomial theorem gives exactly \cref{eq:q-extraction-gaussian}.
This form makes the mechanism transparent: the filter preserves the constant
mode and kills every permitted decaying mode.

\subsection{First universal filters}

For \(Q=1/4\), the first cases are
\begin{align}
 M=0:&\qquad \upf(x)=q_n,\label{eq:filter0}\\
 M=1:&\qquad \upf(x)=\frac{4q_{n+1}-q_n}{3},\label{eq:filter1}\\
 M=2:&\qquad \upf(x)=\frac{q_n-20q_{n+1}+64q_{n+2}}{45},\label{eq:filter2}\\
 M=3:&\qquad \upf(x)=
 \frac{-q_n+84q_{n+1}-1344q_{n+2}+4096q_{n+3}}{2835}.
 \label{eq:filter3}
\end{align}
The same filter applies to every numerator at a given denominator-depth band:
\(M=1\) for depths \(2,3\), \(M=2\) for depths \(4,5\), and so forth.

\section{Recovery of every geometric coefficient}

The limit requires only the first row of the inverse Vandermonde matrix.  For
completeness, this section gives all rows explicitly.

Let
\begin{equation}\label{eq:Dj}
 D_j=\prod_{\substack{0\le\ell\le M\\\ell\ne j}}(Q^j-Q^\ell)
 =(-1)^jQ^{j(j-1)/2+j(M-j)}(Q;Q)_j(Q;Q)_{M-j}.
\end{equation}
Define the elementary symmetric quantity
\begin{equation}\label{eq:erj}
 e_r^{(j)}
 =e_r\bigl(1,Q,\ldots,\widehat{Q^j},\ldots,Q^M\bigr)
 =[t^r]\frac{(-t;Q)_{M+1}}{1+tQ^j},
\end{equation}
where the hat means omission and \([t^r]\) extracts a coefficient.  The finite
\(q\)-binomial theorem also gives
\begin{equation}\label{eq:erj-qbinom}
 e_r^{(j)}=
 \sum_{h=0}^{r}(-1)^hQ^{jh+\binom{r-h}{2}}
 \qbinom{M+1}{r-h}{Q}.
\end{equation}

\begin{proposition}[Closed inverse for all modes]\label{prop:all-coefficients}
With the notation of \cref{eq:scaled-coefficients},
\begin{equation}\label{eq:all-B}
 B_m=\sum_{j=0}^{M}q_{n_0+j}
 \frac{(-1)^{M-m}e_{M-m}^{(j)}}{D_j},
 \qquad 0\le m\le M,
\end{equation}
and hence
\begin{equation}\label{eq:all-C}
 C_m(x)=Q^{-m(n_0+1)}
 \sum_{j=0}^{M}q_{n_0+j}
 \frac{(-1)^{M-m}e_{M-m}^{(j)}}{D_j}.
\end{equation}
For \(m=0\), this reduces to \cref{eq:q-extraction}.
\end{proposition}

\begin{proof}
The Lagrange basis polynomial for the node \(Q^j\) is
\[
 L_j(z)=\frac{\prod_{\ell\ne j}(z-Q^\ell)}{D_j}.
\]
Expanding its numerator in elementary symmetric functions gives
\[
 [z^m]L_j(z)=\frac{(-1)^{M-m}e_{M-m}^{(j)}}{D_j}.
\]
Since \(P(z)=\sum_jq_{n_0+j}L_j(z)\), coefficient comparison proves
\cref{eq:all-B}; the rescaling in \cref{eq:scaled-coefficients} proves
\cref{eq:all-C}.  Formula \cref{eq:erj-qbinom} follows by expanding
\((-t;Q)_{M+1}\) and multiplying by the geometric series for
\((1+tQ^j)^{-1}\), truncated at degree \(r\).
\end{proof}

In practice, exact Gaussian elimination on \cref{eq:q-vandermonde-system} is
often simpler than evaluating \cref{eq:all-C}.  The proposition is useful
conceptually: it proves that every derivative combination
\((-1)^ma_m\upf^{(2m)}(x)\) is a rational expression in finitely many rational
spline values.

\section{An exact \(q\)-binomial recurrence and model diagnostics}

Because \(q_n\) is a linear combination of the modes with eigenvalues
\(1,Q,\ldots,Q^M\), it is annihilated by their characteristic polynomial.

\begin{theorem}[Eventual annihilating recurrence]\label{thm:recurrence}
For every \(n\ge d\),
\begin{equation}\label{eq:recurrence-operator}
 \prod_{m=0}^{M}(\Shift-Q^m)q_n=0.
\end{equation}
Equivalently,
\begin{equation}\label{eq:recurrence-qbinom}
 \boxed{\displaystyle
 \sum_{j=0}^{M+1}(-1)^{M+1-j}
 Q^{\binom{M+1-j}{2}}
 \qbinom{M+1}{j}{Q}\,q_{n+j}=0.}
\end{equation}
\end{theorem}

\begin{proof}
Each factor \(\Shift-Q^m\) kills the corresponding geometric mode in
\cref{eq:main-geometric}.  Expanding
\[
 \prod_{m=0}^{M}(z-Q^m)=z^{M+1}(z^{-1};Q)_{M+1}
\]
by the finite \(q\)-binomial theorem yields \cref{eq:recurrence-qbinom}.
\end{proof}

This recurrence provides three useful diagnostics.
\begin{enumerate}[leftmargin=2.2em]
\item \textbf{Threshold detection.}  If one computes exact prefix values
      without trusting the theoretical threshold, the first level from which
      \cref{eq:recurrence-qbinom} holds on several successive blocks identifies
      the onset of the geometric regime.  The theorem guarantees onset no later
      than \(n=d\).
\item \textbf{Model-order verification.}  Testing orders
      \(0,1,\ldots,M\) reveals whether accidental lower-mode cancellations
      reduce the effective model.  For a reduced nonintegral dyadic point the
      degree is exactly \(M\), so order \(M\) is necessary.
\item \textbf{Error detection.}  In exact rational arithmetic, a nonzero
      residual is an unambiguous implementation error or an indication that
      the starting index is still in the transient range.
\end{enumerate}

\section{Stability of the extraction filter}

Although the weights alternate in sign, they remain uniformly tame.  From
\cref{eq:weight-product},
\begin{equation}\label{eq:weight-norm}
 \sum_{j=0}^{M}|W_{M,j}|
 =\frac{\prod_{r=1}^{M}(1+Q^r)}{\prod_{r=1}^{M}(1-Q^r)}
 =\frac{(-Q;Q)_M}{(Q;Q)_M}.
\end{equation}
For \(Q=1/4\), this increases to the finite limit
\begin{equation}\label{eq:weight-limit}
 \frac{(-Q;Q)_\infty}{(Q;Q)_\infty}
 =1.969260353668\ldots.
\end{equation}
Consequently, if approximate input values have absolute errors at most
\(\varepsilon\), the extracted value has absolute error at most
\[
 1.970\,\varepsilon.
\]
There is no depth-dependent explosion of the extrapolation coefficients.  Of
course, when exact rational prefix values are supplied, the extraction itself
is exact.

The bounded norm is unusually favorable for high-order extrapolation.  It is a
consequence of the strongly separated geometric nodes \(1,1/4,1/16,\ldots\),
not a generic property of polynomial extrapolation.

\section{Worked exact examples}

The following examples use the user's indexing: \(\upf_n\) contains the
factors \(k=0,\ldots,n\).

\subsection{Depth two: \(x=1/4\)}

Here \(d=2\), \(M=1\), and the geometric formula begins at \(n=2\).  Exact
spline evaluation gives
\[
 q_2=\frac{15}{16},
 \qquad
 q_3=\frac{179}{192}.
\]
The universal two-level filter \cref{eq:filter1} yields
\[
 \upf\!\left(\frac14\right)
 =\frac{4q_3-q_2}{3}
 =\frac{67}{72}.
\]
Solving the two-by-two system gives the complete sequence
\begin{equation}\label{eq:example-quarter}
 \boxed{\displaystyle
 \upf_n\!\left(\frac14\right)
 =\frac{67}{72}+\frac49\,4^{-(n+1)},
 \qquad n\ge2.}
\end{equation}
The mode coefficient agrees with the differential formula because
\(a_1=1/18\) and \(\upf''(1/4)=-8\).

\subsection{Depth three: \(x=1/8\)}

Here \(d=3\), but still \(M=1\).  The first guaranteed samples are
\[
 q_3=\frac{383}{384},
 \qquad
 q_4=\frac{1531}{1536}.
\]
Therefore
\[
 \upf\!\left(\frac18\right)
 =\frac{4q_4-q_3}{3}
 =\frac{287}{288},
\]
and the full eventual formula is
\begin{equation}\label{eq:example-eighth}
 \boxed{\displaystyle
 \upf_n\!\left(\frac18\right)
 =\frac{287}{288}+\frac29\,4^{-(n+1)},
 \qquad n\ge3.}
\end{equation}

\subsection{Depth four: \(x=1/16\)}

Now \(d=4\) and \(M=2\), so two decaying modes occur.  The first three exact
values in the guaranteed regime are
\begin{align*}
 q_4&=\frac{24575}{24576},\\
 q_5&=\frac{1965959}{1966080},\\
 q_6&=\frac{94365509}{94371840}.
\end{align*}
Applying \cref{eq:filter2},
\begin{equation}\label{eq:example-sixteenth-limit}
 \upf\!\left(\frac1{16}\right)
 =\frac{q_4-20q_5+64q_6}{45}
 =\frac{2073457}{2073600}.
\end{equation}
Solving the complete Vandermonde system gives
\begin{equation}\label{eq:example-sixteenth}
 \boxed{\displaystyle
 \upf_n\!\left(\frac1{16}\right)
 =\frac{2073457}{2073600}
  +\frac{5}{162}\,4^{-(n+1)}
  -\frac{3968}{2025}\,16^{-(n+1)},
 \qquad n\ge4.}
\end{equation}
From \(a_1=1/18\) and \(a_2=31/16200\), this corresponds to
\[
 \upf''\!\left(\frac1{16}\right)=-\frac59,
 \qquad
 \upf^{(4)}\!\left(\frac1{16}\right)=-1024.
\]

\subsection{Same depth, different numerator: \(x=3/16\)}

The extraction weights are unchanged because they depend only on \(d\).  The
samples
\[
 q_4=\frac{8011}{8192},\qquad
 q_5=\frac{1922041}{1966080},\qquad
 q_6=\frac{92250811}{94371840}
\]
give
\begin{equation}\label{eq:example-three-sixteenth}
 \upf\!\left(\frac3{16}\right)=\frac{2026943}{2073600}
\end{equation}
and
\begin{equation}\label{eq:example-three-sixteenth-sequence}
 \upf_n\!\left(\frac3{16}\right)
 =\frac{2026943}{2073600}
  +\frac{67}{162}\,4^{-(n+1)}
  +\frac{3968}{2025}\,16^{-(n+1)},
 \qquad n\ge4.
\end{equation}
The universal filter is the same; only the mode amplitudes change.

\subsection{Degenerate low-depth cases}

At the three integer points,
\[
 \upf_n(-1)=0,\qquad \upf_n(0)=1,\qquad \upf_n(1)=0
\]
for every \(n\ge0\).  At \(x=\pm1/2\), \(d=1\) and \(M=0\), so
\[
 \upf_n(\pm1/2)=\upf(\pm1/2)=\frac12,
 \qquad n\ge1.
\]
These are exact stabilization rather than nontrivial extrapolation.

\section{A practical exact algorithm}

\subsection{Input and guaranteed sample count}

Given a dyadic rational \(x\in[-1,1]\):
\begin{enumerate}[leftmargin=2.2em]
\item Reduce \(x\) to lowest terms and determine
      \(d=\min\{r:2^rx\in\Z\}\).
\item Set \(M=\lfloor d/2\rfloor\) and \(Q=1/4\).
\item Choose any starting level \(n_0\ge d\); the smallest guaranteed choice is
      \(n_0=d\).
\item Compute the \(M+1\) exact rational values
      \(q_{n_0},q_{n_0+1},\ldots,q_{n_0+M}\), for example from
      \cref{eq:truncated-power}.
\item Return the dot product in \cref{eq:q-extraction-gaussian}.
\end{enumerate}
This returns \(\upf(x)\) exactly.  No limiting process, numerical tolerance, or
rational reconstruction is needed.

\subsection{Recovering and certifying the complete sequence}

To obtain every coefficient, solve \cref{eq:q-vandermonde-system} exactly and
undo the scale:
\[
 C_m=Q^{-m(n_0+1)}B_m.
\]
Then certify the result in either of two ways:
\begin{enumerate}[leftmargin=2.2em]
\item compare the predicted formula with one or more additional exact values
      \(q_{n_0+M+1},q_{n_0+M+2},\ldots\); or
\item evaluate the exact recurrence residual in
      \cref{eq:recurrence-qbinom}.
\end{enumerate}
The proof already guarantees correctness, but these checks are valuable for
software regression testing.

\subsection{Complexity}

The direct truncated-power sum uses \(O(2^{n+1})\) terms.  Since the minimal
sample block reaches roughly \(n=d+\lfloor d/2\rfloor\), this transparent
method is exponential in the dyadic depth.  The extrapolation itself is cheap:
weight construction takes \(O(M^2)\) rational operations if powers and
\(q\)-Pochhammers are accumulated, and a full dense solve takes \(O(M^3)\).
For deep points, the prefix values can be accelerated by recursive
Thue--Morse prefix-power sums, repeated-integration recurrences, or existing
exact dyadic algorithms \cite{AriasArithmetic,Haugland2016}.  The extraction
formula is independent of how the input prefix values are obtained.

\section{Exact-arithmetic experiments}

The archive accompanying this article contains
\begin{center}
\texttt{verify\_dyadic\_extraction.py}
\end{center}
and its captured output.  The script implements:
\begin{enumerate}[leftmargin=2.2em]
\item the exact truncated-power spline formula with Thue--Morse signs;
\item \((Q;Q)_r\), Gaussian \(q\)-binomials, and the extraction weights;
\item exact Gauss--Jordan solution of the \(q\)-Vandermonde system;
\item prediction of unused spline levels;
\item the annihilating recurrence \cref{eq:recurrence-qbinom};
\item an exhaustive regression over every reduced dyadic point of depth at
      most six.
\end{enumerate}
Every calculation uses Python's \texttt{Fraction}; a statement such as
``residual zero'' therefore means exact equality of integers after reduction,
not agreement to a floating-point tolerance.

Representative output reproduces the examples in the preceding section.  The
exhaustive test checks \(129\) points, including the three integers, and finds
that every held-out geometric prediction and every recurrence residual is
exactly zero.  This computation is ancillary: it tests the formulas and their
indexing, while the proofs establish them for all dyadic depths.

\section{Further consequences and extensions}

\subsection{Derivative extraction}

The local theorem contains more than values.  If \(0\le r\le d\), then for
\(n\ge d\),
\begin{equation}\label{eq:derivative-geometric}
 \upf_n^{(r)}(x)=
 \sum_{m=0}^{\lfloor(d-r)/2\rfloor}
 (-1)^ma_m\upf^{(r+2m)}(x)Q^{m(n+1)}.
\end{equation}
Thus the same \(q\)-binomial filter, with
\(M_r=\lfloor(d-r)/2\rfloor\), extracts \(\upf^{(r)}(x)\) exactly from
\(M_r+1\) consecutive derivative approximants.  If \(r>d\), both the true
derivative and the local spline derivative vanish once the centered-cell
regime is reached.

\subsection{Using a larger stride}

Suppose samples are available at levels
\(n,n+s,n+2s,\ldots,n+Ms\).  The corresponding interpolation nodes are
\[
 1,Q^s,Q^{2s},\ldots,Q^{Ms}.
\]
Replacing \(Q\) by \(Q^s\) everywhere in
\cref{eq:q-extraction-gaussian} gives another exact formula.  This can be useful
when only every \(s\)-th prefix has been computed.

\subsection{Arbitrary selected levels}

For distinct offsets \(r_0,\ldots,r_M\), ordinary Lagrange interpolation at
nodes \(Q^{r_j}\) gives
\begin{equation}\label{eq:arbitrary-levels}
 \upf(x)=\sum_{j=0}^{M}q_{n+r_j}
 \prod_{\ell\ne j}\frac{-Q^{r_\ell}}{Q^{r_j}-Q^{r_\ell}}.
\end{equation}
The consecutive case is distinguished because its products collapse to
\(q\)-Pochhammer and Gaussian-binomial expressions.

\subsection{A general local-deconvolution principle}

The proof of \cref{thm:local-deconvolution} is not special to sinc beyond the
coefficient series.  Let \(K\) be any compactly supported even probability
density with moment-generating function \(B_K\), and suppose a target density
has an exact factorization
\[
 f=p_n*K_{\delta_n}.
\]
If the support of \(K_{\delta_n}\), when centered at a point \(x\), lies in one
polynomial cell of \(p_n\), then \(p_n(x)\) is obtained from the finite jet of
\(f\) by the reciprocal series \(1/B_K\).  If that jet terminates and
\(\delta_n^2\) follows a geometric progression, the same exact geometric-mode
and Lagrange-extraction theory follows.  The Rvachev case is exceptionally
clean because the tail is a scaled copy of the whole density and
\(\delta_n^2=4^{-(n+1)}\).

\subsection{Formalization roadmap}

The result is well suited to a modular formal proof.  A possible decomposition
is:
\begin{enumerate}[leftmargin=2.2em]
\item prove the uniform knot formula and the centered-cell lemma;
\item represent convolution with the scaled tail on the finite polynomial jet;
\item prove the reciprocal moment-series inversion in a nilpotent algebra;
\item formalize the Thue--Morse derivative identity and dyadic jet termination;
\item invoke the existing geometric Lagrange/\(q\)-binomial infrastructure to
      derive \cref{eq:q-extraction-gaussian} and
      \cref{eq:recurrence-qbinom}.
\end{enumerate}
The finite-dimensional inversion avoids analytic convergence issues at the
main exactness step.

\section{Conclusion}

For a fixed dyadic rational, the finite-sinc-product approximants do not merely
approach the Rvachev up-function with an all-orders expansion.  After a precise
and finite transient, the expansion terminates exactly.  If \(d\) is the
reduced dyadic depth, then every level \(n\ge d\) satisfies
\[
 \upf_n(x)=\upf(x)+\sum_{m=1}^{\lfloor d/2\rfloor}
 C_m(x)4^{-m(n+1)}.
\]
The reason is the conjunction of two exact geometries: the tail convolution is
confined to one centered polynomial spline cell, and the up-function has a
finite Taylor jet at the dyadic point.

The true rational value is therefore the value at zero of a degree
\(\lfloor d/2\rfloor\) polynomial sampled on a geometric grid.  Its Lagrange
formula is the universal filter
\[
 \upf(x)=\frac{1}{(1/4;1/4)_M}
 \sum_{j=0}^{M}(-1)^{M-j}4^{-\binom{M-j+1}{2}}
 \qbinom{M}{j}{1/4}\,\upf_{n+j}(x),
 \qquad M=\left\lfloor\frac d2\right\rfloor.
\]
This formula uses only finitely many exact rational spline values, proves their
limit is rational, recovers the limit without approximation, and fits directly
into the \(q\)-binomial interpolation framework already developed around the
Fabius--Rvachev system.

\appendix

\section{A compact proof of the weight identities}

For convenience, put
\[
 W_j=(-1)^{M-j}
 \frac{Q^{\binom{M-j+1}{2}}}{(Q;Q)_j(Q;Q)_{M-j}}.
\]
The Lagrange argument proves
\begin{equation}\label{eq:moment-cancel-appendix}
 \sum_{j=0}^{M}W_jQ^{jr}=\delta_{r0},
 \qquad 0\le r\le M.
\end{equation}
One can also derive this directly from the finite \(q\)-binomial theorem.  The
shift polynomial is
\begin{align*}
 \prod_{r=1}^{M}(z-Q^r)
 &=\sum_{j=0}^{M}(-1)^{M-j}
   Q^{\binom{M-j+1}{2}}
   \qbinom{M}{j}{Q}z^j.
\end{align*}
Dividing by its value at \(z=1\), namely \((Q;Q)_M\), gives the extraction
filter.  At \(z=Q^r\), \(1\le r\le M\), it vanishes; at \(z=1\), it equals one.
This is exactly \cref{eq:moment-cancel-appendix}.

\section{Index conversion to the repository convention}

Some source documents denote by \(p_N\) the inverse transform of the first
\(N\) factors \(k=0,\ldots,N-1\).  The present article follows the user's
notation \(\upf_n\) for factors \(k=0,\ldots,n\).  The conversion is
\[
 N=n+1,
 \qquad
 p_N=\upf_n,
 \qquad
 2^{-N}=2^{-(n+1)}=\delta_n.
\]
Thus a repository expansion in powers of \(4^{-N}\) becomes the present
expansion in powers of \(4^{-(n+1)}\).  Keeping this one-level shift explicit is
important when reproducing the rational examples.

\small
\begin{thebibliography}{99}

\bibitem{Rvachev1990}
V.~A. Rvachev,
\newblock Compactly supported solutions of functional-differential equations and their applications,
\newblock \emph{Russian Mathematical Surveys} \textbf{45} (1990), no.~1, 87--120.

\bibitem{Arias1982}
J.~Arias de Reyna,
\newblock Definici\'on y estudio de una funci\'on indefinidamente diferenciable de soporte compacto,
\newblock \emph{Revista de la Real Academia de Ciencias Exactas, F\'isicas y Naturales de Madrid}
\textbf{76} (1982), no.~1, 21--38;
English translation, \emph{An infinitely differentiable function with compact support: Definition and properties},
\href{https://arxiv.org/abs/1702.05442}{arXiv:1702.05442} (2017).

\bibitem{AriasArithmetic}
J.~Arias de Reyna,
\newblock Arithmetic of the Fabius function,
\newblock \emph{INTEGERS} \textbf{18} (2018), Article A51;
\href{https://arxiv.org/abs/1702.06487}{arXiv:1702.06487}.

\bibitem{Haugland2016}
J.~K. Haugland,
\newblock Evaluating the Fabius function,
\newblock \href{https://arxiv.org/abs/1609.07999}{arXiv:1609.07999} (2016).

\bibitem{deBoor}
C.~de Boor,
\newblock \emph{A Practical Guide to Splines}, revised edition,
\newblock Springer, New York, 2001.

\bibitem{GasperRahman}
G.~Gasper and M.~Rahman,
\newblock \emph{Basic Hypergeometric Series}, second edition,
\newblock Encyclopedia of Mathematics and its Applications, Vol.~96,
Cambridge University Press, 2004.

\bibitem{ReshetnikovRepo}
V.~Reshetnikov,
\newblock \emph{ProveIt: Fabius Function documentation and semi-formalized research frontiers},
\newblock \url{https://github.com/VladimirReshetnikov/ProveIt/tree/main/Analysis/FabiusFunction/docs},
accessed 2 September 2026.  In particular, see the finite dyadic sinc-product and
geometric \(q\)-series synthesis under the \texttt{semi-formalized-research-frontiers} tree.

\end{thebibliography}

\end{document}
